\section*{Chapitre \ref{ch:in:crypto-native-firms} --- Sociétés crypto-natives}

\begin{solution}{exo:in:crypto-native-firms:1}
Revenu~: $-59.3~\%$ (2022), $-2.7~\%$ (2023), $+111.2~\%$ (2024), $+9.4~\%$ (2025). Effectifs~: $-24.3~\%$ (2023), $+10.4~\%$
(2024), $+31.3~\%$ (2025)~; les effectifs de 2021 ne figurent pas dans les rapports consultés.
\end{solution}

\begin{solution}{exo:in:crypto-native-firms:2}
25~\% (100\,000~\$ au prix d'attribution) à l'échéance du cliff, puis $1/48$ de l'attribution (8\,333~\$) chaque mois pendant 36 mois.
\end{solution}

\begin{solution}{exo:in:crypto-native-firms:3}
Impôt $0.4\times10\,000=\$4\,000$~; produit 5\,000~\$~; net 1\,000~\$.
\end{solution}

\begin{solution}{exo:in:crypto-native-firms:4}
Sans dérive, le prix espéré reste constant, mais la médiane lognormale décroît comme $e^{-\sigma^2t/2}$~: quelques trajectoires montent fortement et beaucoup baissent.
La moyenne repose sur les rares issues très élevées~; l'issue typique est inférieure à elle.
\end{solution}

\begin{solution}{exo:in:crypto-native-firms:5}
17\,150~\$~; le solde de 42\,850~\$ relève d'une créance chirographaire générale.
\end{solution}

\begin{solution}{exo:in:crypto-native-firms:6}
2023 (effectifs $-24~\%$, le revenu n'ayant baissé que de 3~\%) et 2025 (effectifs $+31~\%$, le revenu ayant progressé de 9~\%)~: dans les deux cas, les effectifs
réagissaient au revenu de l'année précédente, à l'effondrement de 2022 et au doublement de 2024.
\end{solution}

\begin{solution}{exo:in:crypto-native-firms:7}
La médiane passe de 106\,193~\$ à 76\,275~\$, et la part des trajectoires où une tranche a un impôt supérieur à son produit monte
de 76.2~\% à 86.8~\%~: un blocage plus long laisse plus de temps au prix pour baisser après la fixation de l'impôt.
\end{solution}

\begin{solution}{exo:in:crypto-native-firms:8}
Espérance égale ne signifie pas équivalence~: la médiane de l'attribution est inférieure à la moitié de sa valeur en espèces, et ses issues vont de
presque rien à plus du double de la valeur en espèces, et l'impôt porte sur la valeur à l'acquisition, quoi qu'il arrive ensuite. Un salarié réfractaire au risque
valorise l'attribution sous son espérance (équivalent certain du chapitre~13), et la survie de la société elle-même est corrélée au
jeton.
\end{solution}

\begin{solution}{pb:in:crypto-native-firms:1}
\begin{enumerate}
\item Teneurs de marché et sociétés de trading, fonds, bourses et courtiers, équipes de protocole et fondations.
\item Une association à but non lucratif qui finance et coordonne un protocole open source à partir d'une trésorerie partiellement libellée en son jeton~; elle vend périodiquement une partie de son
jeton et constitue une épargne en monnaie légale sur les marchés haussiers, pour financer les dépenses sur les marchés baissiers.
\item Marchés ouverts jour et nuit~; auto-conservation~; \omterm{def:in:crypto-native-firms:tokens}{rémunération en jetons}~; statut réglementaire non tranché.
\item Revenu 7.8, 3.2, 3.1, 6.6, 7.2 milliards de dollars~; salariés (2022--2025) 4\,510, 3\,416, 3\,772, 4\,951.
\item Juin 2022, environ 18~\% (environ 1\,100)~; janvier 2023, environ 21~\% des effectifs de fin 2022 (environ 950).
\item Rémunération en crypto-actif selon un calendrier d'acquisition (vesting), souvent assortie d'un blocage~; un \omterm{def:in:crypto-native-firms:tokens}{bon de souscription de jetons} (token warrant) attaché à des titres de capital pour recevoir de futurs
jetons.
\item Comme salaires (États-Unis) et comme revenus d'emploi soumis à l'impôt sur le revenu et à la National Insurance (Royaume-Uni).
\item En général à l'acquisition, sur la valeur à ce moment-là.
\item Le salarié a payé l'impôt sur une valeur qui a disparu.
\item 1\,000~\$.
\item 25~\% à douze mois, puis $1/48$ par mois jusqu'au mois 48~; chaque tranche bloquée six mois.
\item Calendrier~; simulation des prix~; impôt à l'acquisition et produit après blocage, par tranche~; distribution contre espèces.
\item 106\,193~\$~; 17\,467~\$ et 537\,810~\$~; 235\,062~\$.
\item 240\,000~\$ (400\,000~\$ moins 40~\% d'impôt)~; dans 74.3~\% des trajectoires, l'attribution vaut moins.
\item 76.2~\% à 80~\% de volatilité~; 2.4~\% à 40~\%.
\item Médiane 106\,193~\$, 10\up{e}--90\up{e}~centile de 17\,467~\$ à 537\,810~\$ contre 240\,000~\$ en espèces~; 76.2~\% des trajectoires avec une
tranche imposée sur plus que ce qu'elle rapporte à la vente.
\item Au plus 17\,150~\$.
\item Rejoindre la société au sommet d'un cycle la place dans la cohorte recrutée au pic de revenu et la plus exposée à la prochaine baisse.
\item Vendre les tranches dès que c'est permis~; se couvrir avec des contrats à terme ou des options sur le jeton lorsqu'ils existent et que son contrat
le permet~; négocier des espèces pour une partie de l'attribution, ou payer l'impôt à la vente.
\item L'offre en espèces lui vaut davantage, sauf si elle veut jouer à la loterie sur le jeton et peut supporter d'en perdre l'essentiel~; si elle
prend l'attribution, elle doit prévoir de vendre selon le calendrier et réserver en espèces le montant de l'impôt.
\end{enumerate}
\end{solution}

\begin{solution}{iq:in:crypto-native-firms:1}
La conservation des actifs clients (gestion des clés, portefeuilles), les dépôts et retraits on-chain avec confirmations, un fonctionnement en continu
sans enchère de clôture, et la sécurité face à des attaquants capables de voler de manière irréversible.
\iqlookfor{conservation et irréversibilité.}
\end{solution}

\begin{solution}{iq:in:crypto-native-firms:2}
Les marchés ne ferment jamais et bougent le plus quand les autres dorment~; un desk de trading follow-the-sun transmet le carnet d'ordres d'une région à l'autre, avec une
passation claire des positions, des limites et des incidents ouverts, et un ingénieur d'astreinte par poste.
\iqlookfor{rigueur de la passation.}
\end{solution}

\begin{solution}{iq:in:crypto-native-firms:3}
$e^{-\sigma^2/2}=e^{-0.32}=0.726$ du prix du jour.
\iqlookfor{la médiane lognormale.}
\end{solution}

\begin{solution}{iq:in:crypto-native-firms:4}
Vendre à découvert des contrats à terme ou des contrats perpétuels sur le jeton, ou acheter des puts, s'ils existent et sont liquides~; limites~: le contrat de travail ou les règles
d'initiés peuvent l'interdire, le marché peut être trop étroit et les coûts de financement peuvent être élevés.
\iqlookfor{instruments, et contraintes contractuelles et pratiques.}
\end{solution}

\begin{solution}{iq:in:crypto-native-firms:5}
Flux de retraits attendus et sa queue de distribution, face au risque de garder les clés en ligne~: assez de portefeuilles chauds pour honorer
les retraits, le reste en portefeuilles froids. Risques~: vol depuis les portefeuilles chauds, retards et erreurs opérationnelles lors du transfert depuis le stockage à froid.
\iqlookfor{liquidité contre sécurité.}
\end{solution}

\begin{solution}{iq:in:crypto-native-firms:6}
On corrèle les variations d'effectifs aux variations de revenu à décalage zéro et à décalage d'une période~: ici la relation à décalage d'une période est la plus forte. Avec
quatre ou cinq points, il faut la présenter comme une tendance et non comme un test, et chercher les données d'autres sociétés avant de généraliser.
\iqlookfor{décalages, et prudence avec cinq points.}
\end{solution}
